% Canonical node 006 · C018; current section path 1.1.6.
% Canonical owner: C018
\cnnaNodeHeading{006}{C018}{Canonical BFS/lexicographic provenance birth schedule}
\cnnaNodePlacement{1.1.6}{D2}

\paragraph*{Position and scientific role.}
\texttt{C018} consumes the finite $b$-ary address grammar of \texttt{C003} and
fixes the active deterministic order in which admissible non-root provenance
slots are considered.  It introduces no new free parameter.  The root has
already been created by \texttt{C002}; consequently it is not an element of
the non-root birth schedule.  The node owns an ordering convention only.  It
does not yet assign an event number, a birth time, a response value, a
conductance, a geometry, or a dynamical update.

\paragraph*{Formal order.}
Let $S_b=\{0,\ldots,b-1\}$ and let $S_b^*$ be the C003 address words.  Local
slot ranks are ordered by the ordinary strict order on their values.  At equal
word length, this induces the lexicographic relation $<_{\mathrm{lex}}$.
C018 then defines
\begin{equation}
  a \prec_{\mathrm{BFSlex}} c
  \quad\Longleftrightarrow\quad
  |a|<|c|
  \;\lor\;
  \bigl(|a|=|c|\ \land\ a<_{\mathrm{lex}}c\bigr).
  \label{eq:c018-bfslex}
\end{equation}
This is the shortlex order on provenance words: depth is compared first and
lexicographic order breaks ties.  It is selected as the active CNNA
construction convention; C003 alone does not force this convention against
all other possible linearizations.
% CNNA-EXTREF-BEGIN EXT-USE-C018-SHORTLEX-DIRECT
\cite[Ch.~1, finite words and lexicographic order]{Lothaire1997CombinatoricsWords}.
% CNNA-EXTREF-END EXT-USE-C018-SHORTLEX-DIRECT

For one parent $u$ and ranks $r<s$, the order satisfies
\begin{equation}
  u\mathbin{\Vert}(r) \prec_{\mathrm{BFSlex}}
  u\mathbin{\Vert}(s),
  \label{eq:c018-sibling-order}
\end{equation}
and every parent precedes every child,
\begin{equation}
  u \prec_{\mathrm{BFSlex}} u\mathbin{\Vert}(r).
  \label{eq:c018-parent-child}
\end{equation}

% CNNA-INFOBOX-BEGIN INFO-C018-SHORTLEX-BASEB
\begin{cnnaInfoBox}{Shortlex order and the fixed-depth base-$b$ coordinate}
At depth $d$, an address $a=(a_0,\ldots,a_{d-1})$ has the ordinary
base-$b$ coordinate $\sum_{j=0}^{d-1}a_j b^{d-1-j}$.  Prefixing this
coordinate by the number of words at smaller depths yields the closed index
\[
  1+\sum_{k=1}^{d-1}b^k+\sum_{j=0}^{d-1}a_j b^{d-1-j}
\]
for the non-root shortlex schedule.  Leading zero slots remain part of the
word, so this is a coordinate for the chosen order rather than an
identification of provenance words with integers.
% CNNA-EXTREF-BEGIN EXT-USE-C018-WORDS-INFOBOX
\cite[Ch.~1]{Lothaire1997CombinatoricsWords}.
% CNNA-EXTREF-END EXT-USE-C018-WORDS-INFOBOX
\end{cnnaInfoBox}
% CNNA-INFOBOX-END INFO-C018-SHORTLEX-BASEB

\paragraph*{Strict-total-order theorem.}
The Lean development proves constructively that the address relation in
Eq.~\eqref{eq:c018-bfslex} is irreflexive, transitive, asymmetric, and
trichotomous:
\begin{equation}
  a\prec_{\mathrm{BFSlex}}c
  \;\lor\;
  a=c
  \;\lor\;
  c\prec_{\mathrm{BFSlex}}a.
  \label{eq:c018-trichotomy}
\end{equation}
Hence distinct provenance addresses are comparable in exactly one direction.
The proof first establishes irreflexivity, transitivity and trichotomy for the
lexicographic word relation and then lifts these facts through the depth-first
disjunction defining \texttt{BirthBefore}.  No classical case split is needed.

\paragraph*{Bounded open slots.}
For a C003 grammar with cutoff $L$, the formal slot object consists of a
bounded parent $u$, a local rank $r$, and a proof that $|u|+1\le L$.  Its
selected child is $u\mathbin{\Vert}(r)$.  The slot relation is defined by
comparing these child addresses with $\prec_{\mathrm{BFSlex}}$.  It is
therefore irreflexive, transitive and asymmetric; slots selecting distinct
children are totally comparable.  The word ``open'' at this node means
\emph{admissible potential birth slot}.  Openness relative to an evolving born set and uniqueness of the least
currently open slot require additional state and are owned by C004.  P002
certifies only the static order exported by C018; that six-declaration proof
contract is now kernel-verified with an empty axiom profile.

\paragraph*{Executable realization and cross-layer agreement.}
Python stores one slot as $(u,r,u\mathbin{\Vert}(r))$ and orders it by
\begin{equation}
  K(u,r)=\bigl(|u|,u,r\bigr).
  \label{eq:c018-python-key}
\end{equation}
A breadth-first parent queue emits ranks $0,\ldots,b-1$ in increasing order
and appends each emitted child to the same queue.  Because every child has
length $|u|+1$ and lexicographic comparison of $u\mathbin{\Vert}(r)$ first
compares $u$ and then $r$, the key in Eq.~\eqref{eq:c018-python-key} is
extensionally the same ordering criterion as the Lean child-address relation.
The executable tests reproduce exact schedules for $(b,L)=(2,3)$ and $(3,2)$,
strictness, transitivity, determinism, cutoff zero, and the finite carrier
cardinality.

\paragraph*{Necessity checks and ownership boundary.}
Removing the depth comparison would produce depth-first lexicographic order,
not breadth-first order.  Removing the final rank from the Python key would
leave siblings of one parent unordered.  Treating the root as a birth slot
would duplicate the C002 genesis.  Conversely, the ordering relation alone does not prove that a state-dependent
selector possesses a unique least open slot, nor that the finite schedule
enumerates every bounded address exactly once.  Static order certification is
isolated in \texttt{P002} and is now closed; state-dependent least-open
selection is owned by \texttt{C004}, and finite exhaustivity by \texttt{P004}.
Those downstream obligations do not weaken the verified C018 order theorem.

\paragraph*{Result and downstream handoff.}
C018 yields a fixed, strict total order on provenance addresses and a derived
strict order on admissible slots.  It selects the first slot for
\texttt{C004A}, orders recurrent candidates for \texttt{C004}, supplies the
ordering input later used by \texttt{C010}, \texttt{C019}, \texttt{C042} and
the control \texttt{CTRL005}, and supports the proof gates \texttt{P002},
\texttt{P003} and \texttt{P004}.  These outgoing edges expose dependencies;
they do not close the downstream nodes.

% CNNA-OPEN-PROVENANCE-BEGIN C018
\begin{cnnaInfoBox}{Open-provenance role: Event provenance versus recurrent live dynamics}
C018 orders birth events in an acyclic provenance history.  Later live-state feedback may be recurrent or cyclic without turning the event-provenance relation itself into a cycle.
\end{cnnaInfoBox}
% CNNA-OPEN-PROVENANCE-END C018

\noindent\textbf{Primary code anchors.}
Python:
\path{derivation/code/python/src/cnna/derivation/s01_primitive_response_coupled_finite_provenance/s01_primitive_inputs_grammar_and_event_order/s06_c018__canonical_bfs_lexicographic_provenance_birth_schedule.py},
\texttt{open\_slot\_key}, lines 46--53; \texttt{canonical\_birth\_slots},
lines 61--88; \texttt{CanonicalBirthSchedule}, lines 97--121.  Lean:
\path{derivation/code/lean/core/src/CNNA/Derivation/S01_PrimitiveResponseCoupledFiniteProvenance/S01_PrimitiveInputsGrammarAndEventOrder/S06_C018_CanonicalBfsLexicographicProvenanceBirthSchedule.lean},
order definitions, lines 23--37; address-order proof chain, lines 39--236;
bounded-slot realization, lines 241--325.  The complete declaration and hash
register is given in the Supplement.
